\documentclass[11pt]{article}
\usepackage{fontspec}
\setmainfont{DejaVu Serif}
\usepackage{amsmath,amssymb}
\usepackage[a4paper,margin=1in]{geometry}\usepackage{microtype}\usepackage{parskip}\usepackage[hidelinks]{hyperref}\pagestyle{empty}
\begin{document}
\begin{center}
{\Large\bfseries Từ bối cảnh sức khỏe dọc đến hành động an toàn: hợp đồng quản trị đồng phiên bản với rào chắn an toàn Tầng~1 phi-LLM cho AI y học gia đình}\\[0.7em]
\textbf{Nguyễn Ngọc Thiện, Trịnh Minh Quang}\\
Trường THPT Hai Bà Trưng, Huế, Việt Nam\\
Tác giả liên hệ: \texttt{kakangocthien109@gmail.com}
\end{center}
\textbf{Chuyên đề:} AI và chuyển đổi số trong Y tế và Giáo dục Y khoa

\textbf{Đặt vấn đề.} Trong y học gia đình và chăm sóc sức khỏe ban đầu, hồ sơ bệnh án theo dõi dọc tích lũy lượng lớn dữ liệu qua nhiều đợt khám, bao gồm tiền sử dùng thuốc, bệnh mạn tính và các thông tin cá nhân nhạy cảm. Một hệ thống AI y tế có trạng thái có thể đọc một ảnh chụp hồ sơ bệnh nhân khi dữ liệu và quyền truy cập còn hợp lệ, nhưng chỉ tạo đề xuất cập nhật sau khi bệnh án, trạng thái đồng thuận hoặc chính sách quản trị đã thay đổi (hiện tượng sai lệch thời gian TOCTOU từ đọc đến ghi). Bên cạnh đó, việc nạp toàn bộ bệnh án không giới hạn vào mô hình ngôn ngữ lớn (LLM) gây nguy cơ lộ lọt thông tin sức khỏe nhạy cảm không liên quan (vi phạm nguyên tắc giảm thiểu dữ liệu Zero-PHI) và dễ dẫn đến ảo giác tính toán thời gian (như hạn tái khám sàng lọc mạn tính, lịch tiêm chủng) hoặc bỏ sót tương tác thuốc nghiêm trọng.

\textbf{Phương pháp.} Nghiên cứu trình bày kiến trúc GLHS được triển khai trong CLARA-Care cho quy trình y học gia đình và chăm sóc ban đầu. GLHS thiết lập \textbf{Rào chắn trạng thái hai tầng (Dual-Layer State Barrier)}: (1) \textbf{Rào chắn an toàn Tầng~1 phi-LLM} (Nhân đối soát xác định tại API Gateway bằng Python/PostgreSQL) chịu trách nhiệm tính toán thời gian hai chiều ($\tau_{\text{valid}} \cap \tau_{\text{known}}$), xác định thời hạn tái khám sàng lọc định kỳ và khoảng thời gian ân hạn tiêm chủng theo phác đồ bằng phép toán thời gian tất định (loại trừ ảo giác thời gian của LLM), kiểm tra tương tác thuốc nghiêm trọng (DDI), xác thực đồng thuận động và quyền RBAC, đồng thời đóng gói mã băm Merkle ($\Pi_{\text{Merkle}}$) thành Ảnh chụp trạng thái sức khỏe giới hạn theo tác vụ (THSS) nhằm thực thi nguyên tắc giảm thiểu dữ liệu theo tác vụ; và (2) \textbf{Tầng~2 (Epistemic Arbiter -- LLM)} chỉ tập trung suy luận ngữ nghĩa định tính, được cách ly hoàn toàn khỏi tính toán số học thời gian thô. Đề xuất bền vững mang theo ràng buộc THSS đến bước Chuyển trạng thái có quản trị (GST) để tái thẩm định độ mới trước khi ghi. Đánh giá gồm kiểm thử hợp đồng xác định, ma trận đồng thời PostgreSQL (12 lịch), đoàn hệ tổng hợp 64 đối tượng trên Claude Sonnet 4.6 và Gemini 3.6 Flash High ($1.152$ ô đánh giá lặp lại), và đối chiếu chẩn đoán ghép cặp cấp đối tượng trên đoàn hệ 384 đối tượng chăm sóc ban đầu độc lập.

\textbf{Kết quả.} Đối chiếu cấp đối tượng ($N=384$ đối tượng, 3.456 ô) ghi nhận 70 thắng, 73 thua và 241 hòa ($\hat{\Delta} = -0{,}781\%$, $SE = 3{,}12\%$, kiểm định dấu $p=0{,}867$). Do khoảng tin cậy bootstrap 95\% $[-6{,}77\%, +5{,}21\%]$ cắt qua biên độ hẹp $\pm 2{,}0\%$, cỡ mẫu $N=384$ chưa đủ lực thống kê để khẳng định tương đương hay không thua kém thống kê (đòi hỏi $N \ge 7.500$). Thay vì khẳng định tính không suy giảm quyết định lâm sàng hay tối ưu Pareto, việc giảm thiểu bối cảnh giúp giảm 87{,}4\% token nhắc (trung bình 412 so với 3.280 tokens, $p < 10^{-40}$) và giảm 68{,}2\% độ trễ suy luận trong khi độ lệch chính xác lâm sàng vẫn nằm trong phương sai thực nghiệm. Đồng thời, hệ thống đảm bảo giảm thiểu phơi nhiễm dữ liệu nhạy cảm ngoài tác vụ (0{,}0\% lộ lọt thực thể trong bộ kiểm thử) và ngăn chặn 256/256 lần ghi sai TOCTOU ($p = 1{,}727 \times 10^{-77}$ so với baseline không ràng buộc). Rào chắn an toàn Tầng~1 phi-LLM thực thi tính toán thời hạn lâm sàng tất định. Trong đoàn hệ 64 đối tượng, Strict THSS đạt đúng hoàn toàn trên mọi trục ở 63/64 đối tượng với Claude (98{,}44\%) và 64/64 với Gemini (100\%). Ma trận 12 lịch PostgreSQL không ghi nhận giao dịch ghi bị cấm khi thu hồi đồng thuận, đổi vai trò, đổi chính sách hoặc trôi dạt trạng thái. Duyệt trạng thái hữu hạn bao phủ 21.361 trạng thái và 90.432 bước chuyển đến độ sâu 5 mà không vi phạm 11 bất biến đã định nghĩa.

\textbf{Kết luận.} GLHS chứng minh rào chắn an toàn Tầng~1 phi-LLM kết hợp hợp đồng đồng phiên bản từ công bố dữ liệu đến ghi thay đổi giúp ngăn chặn xung đột trạng thái và giảm thiểu phơi nhiễm dữ liệu trong chăm sóc ban đầu trong khi duy trì chất lượng suy luận lâm sàng trong phạm vi phương sai thực nghiệm. Bằng chứng hiện tại dựa trên đoàn hệ tổng hợp và thực nghiệm phần mềm; việc ứng dụng thực tế đòi hỏi nghiên cứu lâm sàng tiến cứu.

\textbf{Từ khóa:} tin học y học gia đình; AI chăm sóc ban đầu; giảm thiểu dữ liệu; bối cảnh theo tác vụ; rào chắn an toàn phi-LLM; nhất quán trạng thái; quản trị lâm sàng; an toàn đồng thời.
\end{document}
